\section{OpenMP并行预处理优化}
\subsection{快速单线程与并行解析}
原版逐行使用输入流解析整数。快速版将文件顺序读取到连续缓冲区，使用字符级整数解析并保留原有错误检查。OpenMP版按完整行边界划分缓冲区；相邻区间首尾相接，允许空区间，保证每行恰好处理一次。不同线程将端点和边写入私有数组，按各数组长度计算输出位置后合并。相同顶点在不同分块重复出现不引入写冲突，重复编号由后续排序去重处理。划分、分配、解析、合并及临时缓冲区释放均计入成本。

\subsection{编号映射并行化}
原版融合执行端点二分映射和度数累加。优化版预先分配稠密边数组，各线程只读取排序后的编号数组，并写入不同的边位置；随后串行扫描稠密边累加度数，避免顶点度数的原子竞争。单独保留拆分单线程版本，以评估额外扫描本身的代价。度数累加并未在本轮并行化。

\subsection{并行排序与去重}
顶点编号与规范化边先分段排序，再执行多轮两两归并。每次归并根据输出位置划分子任务，通过稳定的co-rank分界定位两侧输入区间，使最后一轮大归并也能并行。去重时，每段比较当前元素与全局前驱，统计保留项，再按计数前缀复制到独立输出区间，避免跨段重复遗漏。排序和归并属于标准并行算法思路，本文不将其作为原创算法。

该方法需要额外线性临时数组。额外主存、分配与同步是优化代价，不能只报告排序主体时间。定向、CSR前缀和与填充、最终邻居段排序保持串行；GPU计数内核及归约流程保持相同。
